\documentclass[12pt]{article}
\usepackage[utf8]{inputenc}
\usepackage[T1]{fontenc}
\usepackage[margin=1in]{geometry}
\usepackage{setspace}
\usepackage{xcolor}
\usepackage[colorlinks=true, urlcolor=blue!60!black, linkcolor=blue!60!black]{hyperref}
\setlength{\parindent}{0em}
\setlength{\parskip}{0.6em}
\renewcommand{\baselinestretch}{1.0}
\begin{document}
\begin{center}
{\LARGE \textit{Abrams v.\ United States}, 250 U.S.\ 616 (1919)\par}
\vspace{0.5em}
{\large Case Brief --- Two-Page Edition\par}
\end{center}

\section{Memory Jogger}
Wartime leaflets urging a munitions strike were punishable because their likely effects showed intent; Holmes's dissent demanding imminent danger became modern law (\href{https://supreme.justia.com/cases/federal/us/250/616/}{Justia}).

\section{Detailed Case Facts}
In 1918 five Russian-born radicals in New York opposed U.S.\ intervention in Russia and threw leaflets from a window calling Wilson a ``coward'' and urging munitions workers to a ``general strike,'' while repudiating German militarism.

\section{Procedural History}
Convicted of Espionage Act conspiracies and sentenced to as much as twenty years, the defendants sought review by writ of error.

\section{Judicial Votes}
Affirmed, 7--2. Clarke wrote for the Court; Holmes dissented, joined by Brandeis.

\section{Holding}
The First Amendment did not protect the leaflets, and intent to hinder the war could be inferred from the foreseeable effects of a munitions strike, whatever the defendants' motive.

\section{Analysis of Opinions}
\textbf{Majority:} \textit{Schenck} and \textit{Frohwerk} settled the First Amendment issue; ``[m]en must be held to have intended'' the likely effects of their acts. \textbf{Holmes:} the statute required that hindering the war be the actor's aim, and only a ``clear and imminent danger'' justifies punishment; a ``silly leaflet by an unknown man'' posed none. Because ``time has upset many fighting faiths,'' the Constitution trusts ``free trade in ideas.'' \textit{Brandenburg} (1969) later adopted his imminence principle.

\section{Examples: Future Applications}
\textbf{Same-Side:} speech coordinated with a designated foreign terrorist organization (\textit{Humanitarian Law Project}). \textbf{Opposite-Side:} urging dockworkers to strike next month. \textbf{Fence-Sitter:} a viral call to shut down a defense plant ``right now.''

\section{Critique}
Holmes never acknowledged how far he had moved from \textit{Schenck}; critics question whether the market metaphor fits an age of viral disinformation.

\section{Key Quotations}
\begin{quote}
``But when men have realized that time has upset many fighting faiths, they may come to believe even more than they believe the very foundations of their own conduct that the ultimate good desired is better reached by free trade in ideas -- that the best test of truth is the power of the thought to get itself accepted in the competition of the market, and that truth is the only ground upon which their wishes safely can be carried out. That, at any rate, is the theory of our Constitution. It is an experiment, as all life is an experiment.''
\end{quote}
\noindent --- Justice Holmes, dissenting, 250 U.S.\ at 630.

\end{document}
